\documentclass[11pt]{article}

\usepackage[english]{babel}
\usepackage[T1]{fontenc}
\usepackage{amsmath,amssymb}
\usepackage[margin=1in]{geometry}
\usepackage{lmodern}
\usepackage{microtype}
\usepackage{xcolor}
\usepackage[colorlinks=true,citecolor=blue,urlcolor=gray,linkcolor=blue]{hyperref}

\newcommand{\mbf}[1]{\mathbf{#1}}
\newcommand{\bs}[1]{\boldsymbol{#1}}
\newcommand{\dev}{\operatorname{dev}}
\newcommand{\tr}{\operatorname{tr}}

\title{Apparent non-associative plasticity from distinct isochoric and
volumetric mechanisms}
\author{John T. Foster}
\date{August 18, 2026}

\begin{document}

\maketitle

\begin{abstract}
	Pressure-dependent plasticity commonly represents friction and dilatancy with
	a yield function and a distinct plastic potential.  We examine a framework
	with distinct isochoric and volumetric mechanisms.  The finite-deformation
	kinematics follow from \(\mbf F=\mbf A\bar{\mbf F}\), where \(\mbf A\) is
	the distension gradient and \(\bar{\mbf F}\) is the true-deformation
	gradient.  An elastic--plastic split of each factor assigns isochoric
	rearrangement of the solid material to true-plastic deformation and plastic
	change in the volume assigned to it to plastic distension.  An independent
	distension mechanism requires an additional microstructural closure; under
	the Flory specialization, it becomes the conventional plastic volume ratio.
	The isothermal free-energy inequality shows which stress drives each plastic
	rate.  With separate yield conditions, each mechanism obeys an
	associated flow rule.  Their combined plastic strain rate, however, need not
	be normal to either yield surface.

	For isotropic, cohesionless, monotone loading under compression, an exact
	pressure-dependent specialization follows from a stress-dependent dissipation
	potential and a dilatancy constraint.  Its small-strain limit recovers the
	Drucker--Prager yield surface and a distinct plastic potential, with the
	difference between the friction and dilation slopes supplied by dissipative
	shear resistance.  Under the same proportional compressive loading, the
	independent model stays at the stress point where both yield conditions are
	active and accommodates the volume reduction by compaction.  The coupled
	model instead moves along the Drucker--Prager surface while dilating.  The
	construction carries to finite deformation under the Flory closure and
	preserves nonnegative plastic dissipation.  The contribution is a
	mechanism-resolved finite-deformation interpretation of this construction,
	including its closure requirements and the consequences of independent or
	shared activation.  Thus an associated
	evolution law for the two underlying mechanisms becomes the familiar
	non-associated flow rule when written only in terms of pressure and shear
	stress.
\end{abstract}

\section{Introduction}
\label{sec:introduction}

Frictional plasticity models commonly distinguish the surface that activates
plastic loading from the potential that determines the plastic flow
direction.  This distinction permits independent calibration of friction and
dilation angles, but the resulting flow rule is not normal to a single yield
surface in the usual stress variables.  Classical associated plasticity
instead takes
the plastic rate to be normal to the yield surface.  Drucker derived this
structure from work-hardening and ideal-plasticity assumptions, and Bland
showed its equivalence to a plastic potential identical to the yield function
\cite{drucker1950,bland1957}.  Thus conventional frictional models use one
function to activate plastic loading and another to prescribe the observed
dilation rate, whereas classical associated plasticity requires a single
function to serve both roles.

Drucker and Prager introduced the pressure-dependent yield condition used here
as the cohesionless reference model \cite{druckerprager1952}.  At finite strain,
the multiplicative elastic--plastic split and its intermediate configuration
are classical \cite{lee1969}, and maximum-plastic-dissipation formulations
already connect that split, Mandel stress, and associated flow
\cite{simo1988}.  Finite-strain pressure-sensitive models also combine
Mandel-type stresses with inelastic volume change and cap-type response
\cite{bennett2019}.  The factorization used below must therefore be compared
with this established volumetric--isochoric structure rather than treated as a
new multiplicative kinematics.

Thermodynamic formulations track each internal rate together with the force
that does work through it.  Collins and Houlsby showed that an associated flow
rule written in this larger set of forces can become non-associated after the
internal forces are expressed through the usual stresses
\cite[Sec.~5]{collins1997}.  They also gave a
stress-dependent dissipation construction for Drucker--Prager plasticity
\cite[Sec.~5]{collins1997}.  Other convex, energy-based dissipation models
also derive pressure-dependent non-associated flow
\cite{srinivasa2010,ulloa2021}.  Independently, Drumheller's mixture kinematics
split deformation into true deformation and distension, with elastic--plastic
factors and distinct plastic powers for the two mechanisms
\cite[Secs.~3 and 8]{drumheller2000}.  Combining these results assigns specific
finite-deformation carriers to the pressure-dependent plastic rates: shear
rearrangement changes the true deformation, whereas compaction or dilation
changes distension.

This paper uses Drumheller's terminology to identify explicit physical
mechanisms within the Collins--Houlsby thermodynamic construction at finite
deformation.  For a single solid phase, the split
\(\mbf F=\mbf A\bar{\mbf F}\) assigns true-plastic deformation and plastic
distension separate elastic--plastic factors and separate power-conjugate
forces.  In the independent model, each mechanism has its own yield condition
and plastic multiplier.  In the coupled model, a dilatancy constraint and a
stress-dependent dissipation potential give both mechanisms one multiplier.
This coupling produces the Drucker--Prager yield surface and a different
plastic-potential slope.  For the closed finite-deformation model,
the Flory constraint makes the resulting multiplicative split algebraically
equivalent to the classical volumetric--isochoric split.  The small-strain
dissipation construction is the Collins--Houlsby construction; the contribution
developed here is its mechanism-resolved finite-deformation interpretation, the
closure required by the distension variable, and the distinction between
independent and shared activation.  This structure explains how associated
flow for the mechanisms appears non-associated when written in terms of
pressure and shear stress.

Section~\ref{sec:kinematics} defines the finite-deformation split and its
closure, and Section~\ref{sec:thermodynamics} derives the conjugate forces and
plastic power.  Section~\ref{sec:two_mechanisms} examines the associated
two-mechanism theory and its combined plastic strain rate.
Section~\ref{sec:drucker_prager}
constructs the exact small-strain pressure-dependent specialization, which
Section~\ref{sec:finite_lift} extends to finite deformation.

\section{Single-phase kinematics}
\label{sec:kinematics}

Let \(\bs\varphi(\mbf X,t)\) map a material point \(\mbf X\) in the reference
configuration to its current position \(\mbf x\).  The deformation gradient
and its Jacobian are
\begin{equation*}
	\mbf F=\frac{\partial\bs\varphi}{\partial\mbf X},
	\qquad
	J=\det\mbf F>0.
\end{equation*}
The deformation is decomposed into distension and true deformation,
\begin{equation}
	\mbf F=\mbf A\bar{\mbf F},
	\qquad
	J=a\bar J,
	\qquad
	a=\det\mbf A,
	\qquad
	\bar J=\det\bar{\mbf F}.
	\label{eq:distension_split}
\end{equation}
The tensor \(\bar{\mbf F}\) describes deformation of the solid material,
whereas \(\mbf A\) accounts for the surrounding volume assigned to it.
This split is the single-phase specialization of the true-deformation and
distension kinematics introduced for immiscible mixtures by Drumheller
\cite[Secs.~3 and 8.2]{drumheller2000}.

Each factor has an elastic--plastic split,
\begin{subequations}
	\label{eq:plastic_splits}
	\begin{align}
		\mbf A       & =\mbf A^e\mbf A^p,
		\label{eq:distension_plastic_split}               \\
		\bar{\mbf F} & =\bar{\mbf F}^{e}\bar{\mbf F}^{p},
		\qquad
		\det\bar{\mbf F}^{p}=1.
		\label{eq:true_plastic_split}
	\end{align}
\end{subequations}
The isochoric condition assigns all plastic volume change to \(\mbf A^p\).
Drumheller's finite-deformation construction makes the same distinction: the
true-plastic factor is incompressible, whereas the plastic-distension factor
may change volume
\cite[eqs.~(71)--(74)]{drumheller2000}.

The pressure-dependent specialization restricts both elastic and plastic
distension to be isotropic,
\begin{equation*}
	\begin{aligned}
		\mbf A & =a^{1/3}\mbf I,
		       & \mbf A^e        & =(a^e)^{1/3}\mbf I,
		       & \mbf A^p        & =(a^p)^{1/3}\mbf I, \\
		a^e    & =\det\mbf A^e,
		       & a^p             & =\det\mbf A^p,
		       & a               & =a^e a^p,           \\
		\zeta  & =\ln a,
		       & \zeta           & =\zeta^e+\zeta^p.
	\end{aligned}
\end{equation*}
The velocity gradient consequently separates into the spherical distension
rate and the true-deformation rate,
\begin{equation}
	\mbf L=\dot{\mbf F}\mbf F^{-1}
	=\frac{\dot\zeta}{3}\mbf I
	+\dot{\bar{\mbf F}}\bar{\mbf F}^{-1}.
	\label{eq:velocity_gradient}
\end{equation}
The true-plastic velocity gradient is
\begin{equation*}
	\bar{\mbf L}^{p}
	=\dot{\bar{\mbf F}}^{p}(\bar{\mbf F}^{p})^{-1},
	\qquad
	\tr\bar{\mbf L}^{p}=0.
\end{equation*}
Differentiating \eqref{eq:true_plastic_split} separates the true-deformation
rate into its elastic and plastic parts,
\begin{equation}
	\dot{\bar{\mbf F}}\bar{\mbf F}^{-1}
	=\dot{\bar{\mbf F}}^{e}(\bar{\mbf F}^{e})^{-1}
	+\bar{\mbf F}^{e}\bar{\mbf L}^{p}
	(\bar{\mbf F}^{e})^{-1}.
	\label{eq:true_rate_split}
\end{equation}

\subsection{Closure of the kinematic split}
\label{sec:kinematic_closure}

The split in \eqref{eq:distension_split} is not determined by
\(\mbf F\) alone.  Even for isotropic distension, \(J=a\bar J\) fixes only the
product of the distension and true-deformation Jacobians.  The scalar \(a\)
must therefore be supplied by an evolution equation or a kinematic constraint;
once \(a\) is known, \(\mbf A=a^{1/3}\mbf I\) and
\(\bar{\mbf F}=a^{-1/3}\mbf F\).  Its exact kinematic rate is
\begin{equation*}
	\frac{\dot a}{a}
	=\tr\left(\dot{\mbf A}\mbf A^{-1}\right).
\end{equation*}

In Drumheller's mixture theory, constituent mass balance relates this rate to
the rates of volume fraction and mass exchange.  In the absence of mass
exchange and for unit reference volume fraction, the distension is the inverse
of the current volume fraction
\cite[eq.~(36)]{drumheller2000}.  That relation provides
a physical evolution law when the corresponding mixture variables are
retained.  The present single-phase theory does not carry volume fraction or
mass exchange, so an application must instead prescribe a constitutive
evolution or constraint for \(a\).  The flow rule
\eqref{eq:finite_distension_flow} later supplies the plastic rate
\(\dot a^p/a^p\), but it does not by itself determine the total distension.

A standard Flory split into volumetric and shape-changing parts closes the model
\cite{flory1961}.  Foster et al. obtain the same specialization from the
distension split for incompressible solid phases
\cite[Sec.~2, eq.~(17)]{foster2026reactingporous}:
\begin{subequations}
	\label{eq:flory_closure}
	\begin{align}
		\mbf A       & =J^{1/3}\mbf I,
		             & a                & =J,
		\label{eq:flory_distension}           \\
		\bar{\mbf F} & =J^{-1/3}\mbf F,
		             & \bar J           & =1.
		\label{eq:flory_true_deformation}
	\end{align}
\end{subequations}
This choice determines both factors directly from \(\mbf F\).  Together with
\eqref{eq:plastic_splits}, it also gives
\(\det\bar{\mbf F}^{e}=\det\bar{\mbf F}^{p}=1\).  The deviatoric and
volumetric plastic rates remain distinct, and each retains its own contribution
to plastic power.  The interpretation is narrower:
\(a^p\) is the conventional plastic volume ratio rather than an independent
distension that can differ from the total volume change.

The closed case also makes contact with conventional multiplicative
plasticity.  Because the spherical factors commute with the true-deformation
factors, define
\begin{equation*}
	\mbf F^e=(a^e)^{1/3}\bar{\mbf F}^{e},
	\qquad
	\mbf F^p=(a^p)^{1/3}\bar{\mbf F}^{p},
	\qquad
	\mbf F=\mbf F^e\mbf F^p,
	\qquad
	\det\mbf F^p=a^p.
\end{equation*}
Thus the Flory-closed kinematics do not introduce an additional single-phase
degree of freedom.  Their role is to identify the isochoric and volumetric
parts of conventional plastic flow with true-plastic deformation and plastic
distension.  An independently evolving distension is instead appropriate when
a volume fraction, porosity, or another microstructural variable supplies the
additional closure.  Such a model must supply its own evolution law for that
variable; the closed model developed below does not need one.

\section{Free-energy restriction and plastic power}
\label{sec:thermodynamics}

Let \(\psi\) denote the isothermal Helmholtz free energy per unit reference
volume.  For the preceding kinematics and the hardening variables
\(\kappa_d\) and \(\kappa_v\), the model takes
\begin{equation*}
	\psi=\psi
	\left(
	\bar{\mbf F}^{e},\zeta^e,\kappa_d,\kappa_v
	\right).
\end{equation*}
Material frame indifference requires the free energy to remain unchanged when
\(\bar{\mbf F}^{e}\) is subjected to a rigid rotation, and angular-momentum balance
makes \(\bs\tau\) symmetric.  The pressure-dependent specialization below
further assumes an isotropic elastic response, for which the Mandel-type
stress is symmetric.  The local free-energy inequality is
\begin{equation*}
	\mathcal D
	=\mbf P:\dot{\mbf F}-\dot\psi
	\ge 0,
\end{equation*}
where \(\mbf P\) is the first Piola stress.  With the Kirchhoff stress
\(\bs\tau=\mbf P\mbf F^T\), equations \eqref{eq:velocity_gradient} and
\eqref{eq:true_rate_split} separate the mechanical power into distension,
elastic true-deformation, and true-plastic contributions,
\begin{equation*}
	\mbf P:\dot{\mbf F}
	=\left[\bs\tau(\bar{\mbf F}^{e})^{-T}\right]:\dot{\bar{\mbf F}}^{e}
	+\left[(\bar{\mbf F}^{e})^T\bs\tau(\bar{\mbf F}^{e})^{-T}\right]:\bar{\mbf L}^{p}
	+\frac{1}{3}\tr(\bs\tau)\dot\zeta.
\end{equation*}
The corresponding free-energy rate is
\begin{equation*}
	\dot\psi=
	\frac{\partial\psi}{\partial\bar{\mbf F}^{e}}
	:\dot{\bar{\mbf F}}^{e}
	+\frac{\partial\psi}{\partial\zeta^e}
	\left(\dot\zeta-\dot\zeta^p\right)
	+\frac{\partial\psi}{\partial\kappa_d}\dot\kappa_d
	+\frac{\partial\psi}{\partial\kappa_v}\dot\kappa_v.
\end{equation*}
Substitution of the mechanical power and free-energy rate gives the collected
free-energy inequality,
\begin{align}
	\mathcal D={} &
	\left[
		\bs\tau(\bar{\mbf F}^{e})^{-T}
		-\frac{\partial\psi}{\partial\bar{\mbf F}^{e}}
		\right]:\dot{\bar{\mbf F}}^{e}
	+\left[
		\frac{1}{3}\tr\bs\tau
		-\frac{\partial\psi}{\partial\zeta^e}
		\right]\dot\zeta
	\notag                  \\
	              & +\left[
		(\bar{\mbf F}^{e})^T\bs\tau(\bar{\mbf F}^{e})^{-T}
		\right]:\bar{\mbf L}^{p}
	+\frac{\partial\psi}{\partial\zeta^e}\dot\zeta^p
	-\frac{\partial\psi}{\partial\kappa_d}\dot\kappa_d
	-\frac{\partial\psi}{\partial\kappa_v}\dot\kappa_v
	\ge0.
	\label{eq:collected_free_energy_inequality}
\end{align}
For the general distension split, we first take \(a\) as an independent
kinematic variable.  Its scalar rate and the elastic true-deformation rate can
therefore vary independently.  Their coefficients
in \eqref{eq:collected_free_energy_inequality} therefore vanish, giving the
reversible restrictions
\begin{subequations}
	\label{eq:reversible_restrictions}
	\begin{align}
		\frac{\partial\psi}{\partial\bar{\mbf F}^{e}}
		 & =\bs\tau(\bar{\mbf F}^{e})^{-T},
		\label{eq:true_stress_restriction}  \\
		\frac{\partial\psi}{\partial\zeta^e}
		 & =\frac{1}{3}\tr\bs\tau.
		\label{eq:distension_stress_restriction}
	\end{align}
\end{subequations}
Under the Flory closure \eqref{eq:flory_closure}, the admissible elastic
process directions instead separate into the scalar volumetric rate and
isochoric variations satisfying
\(\tr[\dot{\bar{\mbf F}}^e(\bar{\mbf F}^e)^{-1}]=0\).  The resulting
restriction on the deviatoric stress response is
\begin{equation*}
	\dev\left[
		\left(
		\bs\tau(\bar{\mbf F}^{e})^{-T}
		-\frac{\partial\psi}{\partial\bar{\mbf F}^{e}}
		\right)(\bar{\mbf F}^{e})^T
		\right]=\mbf0.
\end{equation*}
Equation~\eqref{eq:distension_stress_restriction} determines the spherical
response.  These are the restrictions used by the closed finite-deformation
specialization.  The residual plastic power is unchanged because
\(\bar{\mbf L}^{p}\) is traceless.

The compression-positive mean Kirchhoff stress, deviatoric true-plastic
driving stress, and hardening forces are
\begin{align*}
	p      & =-\frac{1}{3}\tr\bs\tau,
	\\
	\mbf M & =\dev\left[
		(\bar{\mbf F}^{e})^T\bs\tau
		(\bar{\mbf F}^{e})^{-T}
		\right],
	\\
	Q_d    & =-\frac{\partial\psi}{\partial\kappa_d},
	\qquad
	Q_v=-\frac{\partial\psi}{\partial\kappa_v}.
\end{align*}
Substituting \eqref{eq:reversible_restrictions} into
\eqref{eq:collected_free_energy_inequality} leaves the plastic dissipation,
\begin{equation}
	\mathcal D^p
	=\mbf M:\bar{\mbf L}^{p}
	-p\dot\zeta^p
	+Q_d\dot\kappa_d+Q_v\dot\kappa_v
	\ge0.
	\label{eq:plastic_dissipation}
\end{equation}
Equation~\eqref{eq:plastic_dissipation} shows how each internal mechanism
contributes to plastic power.  The deviatoric stress \(\mbf M\) is conjugate
to the isochoric rate \(\bar{\mbf L}^{p}\), whereas \(-p\) is conjugate to the
logarithmic plastic-distension rate
\(\dot\zeta^p=\dot a^p/a^p\).  These powers are the single-phase counterparts
of the true-plastic and plastic-distension powers in Drumheller's constitutive
theory
\cite[eqs.~(108)--(112)]{drumheller2000}.

\section{Independent mechanisms and their combined response}
\label{sec:two_mechanisms}

The analysis first considers two independent rate-independent mechanisms.  Their
equivalent true-plastic stress \(q\) and equivalent true-plastic strain rate
\(\dot\gamma\) are
\begin{equation}
	q=\sqrt{\frac{3}{2}}\lVert\mbf M\rVert,
	\qquad
	\dot\gamma=\sqrt{\frac{2}{3}}
	\lVert\bar{\mbf L}^{p}\rVert.
	\label{eq:equivalent_quantities}
\end{equation}
The active true-plastic flow developed here is restricted to \(q>0\), for
which the normal is unique and
\begin{equation}
	\mbf N=\frac{3}{2}\frac{\mbf M}{q},
	\qquad
	\bar{\mbf L}^{p}=\dot\gamma\mbf N,
	\qquad
	\mbf M:\bar{\mbf L}^{p}=q\dot\gamma.
	\label{eq:deviatoric_normal}
\end{equation}
The tensor \(\mbf N\) is the outward normal to a level surface of \(q\), and
its normalization makes \(\dot\gamma\) agree with
\eqref{eq:equivalent_quantities}.  The isotropic specialization takes
\(\mbf N\) to be symmetric, which sets the plastic spin to zero in the
intermediate configuration.

True-plastic deformation and distension may yield independently.  The elastic
range is defined by the two conditions
\begin{align*}
	f_d(q,\kappa_d)  & =q-Y(\kappa_d)\le0,
	\\
	f_v(-p,\kappa_v) & =|p|-K(\kappa_v)\le0,
\end{align*}
where \(Y\ge0\) and \(K\ge0\) are the true-plastic and distension yield
strengths, respectively.  On smooth portions of these surfaces, the model
postulates the associated flow rules and hardening laws for the two mechanisms
as
\begin{subequations}
	\label{eq:product_flow_rules}
	\begin{align}
		\bar{\mbf L}^{p}
		 & =\dot\lambda_d\frac{\partial f_d}{\partial\mbf M},
		\qquad
		\dot\kappa_d=\dot\lambda_d h_d,
		\label{eq:true_associated_flow}                       \\
		\dot\zeta^p
		 & =\dot\lambda_v\frac{\partial f_v}{\partial(-p)},
		\qquad
		\dot\kappa_v=\dot\lambda_v h_v,
		\label{eq:distension_associated_flow}
	\end{align}
\end{subequations}
with \(h_d,h_v\ge0\).  Each plastic multiplier satisfies the Kuhn--Tucker
conditions
\begin{equation*}
	\dot\lambda_i\ge0,
	\qquad
	f_i\le0,
	\qquad
	\dot\lambda_i f_i=0,
	\qquad
	i\in\{d,v\}.
\end{equation*}
Active loading also requires consistency,
\begin{equation}
	\dot f_i=0
	\quad\text{when}\quad
	\dot\lambda_i>0,
	\qquad
	i\in\{d,v\}.
	\label{eq:consistency}
\end{equation}
Together with the stress update and the hardening laws in
\eqref{eq:product_flow_rules}, these conditions determine the active
multipliers.

Energetic hardening contributes directly to the dissipation inequality.  The
active conditions and \eqref{eq:product_flow_rules} give
\begin{subequations}
	\label{eq:dissipative_hardening_specialization}
	\begin{gather}
		\mathcal D^p
		=\dot\lambda_d\left(Y+Q_dh_d\right)
		+\dot\lambda_v\left(K+Q_vh_v\right)\ge0,
		\label{eq:hardening_dissipation}                    \\
		Y+Q_dh_d
		\ge0,
		\label{eq:true_hardening_restriction}               \\
		K+Q_vh_v
		\ge0.
		\label{eq:distension_hardening_restriction}
	\end{gather}
\end{subequations}
The restrictions \eqref{eq:true_hardening_restriction} and
\eqref{eq:distension_hardening_restriction} are sufficient for
\eqref{eq:hardening_dissipation} because the mechanisms activate independently.
A purely dissipative hardening specialization takes
\(Q_d=Q_v=0\) whenever \(h_d,h_v>0\), so the hardening variables change the
yield strengths without storing energy.  It reduces the preceding result to
\begin{equation*}
	\mathcal D^p
	=Y\dot\lambda_d+K\dot\lambda_v
	\ge0.
\end{equation*}
More general energetic hardening is admissible when
\eqref{eq:true_hardening_restriction} and
\eqref{eq:distension_hardening_restriction} hold.

The two independent mechanisms can appear non-associated when they are viewed
only through the total plastic strain rate.  At small strain,
\eqref{eq:velocity_gradient} separates this rate into a traceless true-plastic part and a spherical
plastic-distension part.  Let \(\dot{\bs\varepsilon}^{p}_{d}\) denote the
former.  Since \(\dot\zeta^p\) is the plastic volumetric strain rate, its
spherical contribution is \(\dot\zeta^p\mbf I/3\), and hence
\begin{equation}
	\dot{\bs\varepsilon}^{p}
	=\dot{\bs\varepsilon}^{p}_{d}
	+\frac{1}{3}\dot\zeta^p\mbf I.
	\label{eq:observable_plastic_rate}
\end{equation}
When both mechanisms are active, the associated flow rules
\eqref{eq:product_flow_rules} give
\begin{equation*}
	\dot{\bs\varepsilon}^{p}
	=\dot\lambda_d\frac{\partial f_d}{\partial\mbf M}
	+\frac{\dot\lambda_v}{3}
	\frac{\partial f_v}{\partial(-p)}\mbf I.
\end{equation*}
The second term adds volume change to the isochoric flow.  This added part does
not point normal to the shear yield surface \(f_d=0\).  A model that keeps this
surface but uses the total rate \eqref{eq:observable_plastic_rate} therefore
appears non-associated.

Two independent yield conditions generally cannot be replaced by one
non-associated yield surface.  When both mechanisms are active,
\(f_d=0\) and \(f_v=0\) impose two equations, whereas a regular one-surface
model imposes only one.  In addition,
\eqref{eq:consistency} determines \(\dot\lambda_d\) and \(\dot\lambda_v\)
through separate consistency equations without generally fixing the multiplier
ratio required by a chosen plastic potential.  Exact equivalence therefore
requires a shared activation condition or an explicit constitutive coupling.

The possible loading cases show the difference without requiring a particular
elastic law.  At \(f_d=0\) and \(f_v<0\), the independent model permits
isochoric true-plastic flow with \(\dot\zeta^p=0\); at \(f_v=0\) and
\(f_d<0\), it permits purely volumetric plastic flow.  At the intersection,
the two consistency equations determine separate multipliers and do not fix a
universal ratio.  The coupled model below replaces these three possibilities
with one active surface and the fixed ratio
\(\dot\zeta^p/\dot\gamma=\beta\).  Coupling therefore fixes how much volume
change accompanies a given amount of isochoric plastic flow.

The plastic power imposes a further restriction.  Under compression \(p>0\),
dilation has \(\dot\zeta^p>0\), and its contribution
\(-p\dot\zeta^p\) in \eqref{eq:plastic_dissipation} is negative.  Requiring
the distension term alone to be nonnegative permits compaction under
compression but rules out dilation.  A dilating frictional material instead
uses part of the positive isochoric plastic power to offset the negative
volumetric contribution.

\section{Coupled mechanisms and Drucker--Prager plasticity}
\label{sec:drucker_prager}

At small strain, the finite-deformation quantities reduce to
\begin{equation}
	\bs\tau=J\bs\sigma\longrightarrow\bs\sigma,
	\qquad
	\mbf M\longrightarrow\dev\bs\sigma,
	\qquad
	\operatorname{sym}\bar{\mbf L}^{p}
	\longrightarrow\dot{\bs\varepsilon}^{p}_{d},
	\qquad
	\dot\zeta^p\longrightarrow\tr\dot{\bs\varepsilon}^{p}.
	\label{eq:small_strain_limit}
\end{equation}
For aligned flow with \(q>0\), equation~\eqref{eq:deviatoric_normal} gives
\(\mbf M:\bar{\mbf L}^p=q\dot\gamma\).  Neglecting energetic hardening, the
limits in \eqref{eq:small_strain_limit} therefore reduce
\eqref{eq:plastic_dissipation} to the combined plastic power
\begin{equation}
	\mathcal D^p=q\dot\gamma-p\dot\zeta^p.
	\label{eq:small_strain_power}
\end{equation}
The specialization adopts compression as positive and restricts the response to
\(p\ge0\) and \(\dot\gamma\ge0\).  Dilatancy then supplies the kinematic constraint
\begin{equation*}
	\dot\zeta^p=\beta\dot\gamma,
	\qquad
	\beta\ge0.
\end{equation*}
Here, \(\beta\) is the dilation slope.  The following dissipation potential
specifies both the allowed rate pairs and the energy they dissipate:
\begin{equation}
	\mathcal R(p;\dot\gamma,\dot\zeta^p)
	=
	\begin{cases}
		\mu p\dot\gamma,
		         & \dot\gamma\ge0\ \text{and}\
		\dot\zeta^p=\beta\dot\gamma,           \\
		+\infty, & \text{otherwise},
	\end{cases}
	\qquad \mu\ge0,
	\label{eq:dissipation_potential}
\end{equation}
The potential is convex and positively homogeneous.  It includes zero plastic
flow and assigns infinite cost to rate pairs that violate the dilation rule.
This construction adapts the
stress-dependent dissipation and dilation constraint used by Collins and
Houlsby in their Drucker--Prager example
\cite[eqs.~(5.8)--(5.14)]{collins1997}.

To determine the elastic range, hold \(p\) fixed.  Let \(\chi_q\) and \(\chi_v\)
denote the generalized thermodynamic forces conjugate to \(\dot\gamma\) and
\(\dot\zeta^p\), respectively.  Because the conjugate rates are strain-like,
these scalar forces have units of stress and are the reduced counterparts of
the generalized stresses used by Collins and Houlsby
\cite[eqs.~(2.10)--(2.11) and (5.11)--(5.13)]{collins1997}.
At zero plastic rate, admissible forces cannot perform more work through any
rate pair than the dissipation potential assigns to that pair.  Therefore,
\begin{align}
	\chi_q\dot\gamma+\chi_v\dot\zeta^p
	 & \le\mathcal R(p;\dot\gamma,\dot\zeta^p)
	\quad\text{for every rate pair}
	\notag                                     \\
	 & \iff
	\chi_q+\beta\chi_v-\mu p\le0.
	\label{eq:zero_rate_force_condition}
\end{align}
Thus, for each fixed value of \(p\), the admissible driving forces satisfy
\begin{equation}
	f_g(\chi_q,\chi_v;p)
	=\chi_q+\beta\chi_v-\mu p
	\le0.
	\label{eq:generalized_yield}
\end{equation}
On the boundary of this region, associated flow gives the plastic rates and
loading conditions,
\begin{equation}
	\begin{gathered}
		(\dot\gamma,\dot\zeta^p)
		=\dot\lambda\frac{\partial f_g}{\partial(\chi_q,\chi_v)}
		=\dot\lambda(1,\beta), \\
		\dot\lambda\ge0,
		\qquad
		f_g\le0,
		\qquad
		\dot\lambda f_g=0,
		\qquad
		\dot f_g=0\ \text{when}\ \dot\lambda>0.
	\end{gathered}
	\label{eq:generalized_loading_conditions}
\end{equation}
The coupled rate is therefore normal to the boundary in the reduced
generalized-force space \((\chi_q,\chi_v)\) when \(p\) is held fixed.

For the cohesionless model without kinematic hardening, the conjugate forces
coincide with those in the physical power \eqref{eq:small_strain_power},
\begin{equation}
	\chi_q=q,
	\qquad
	\chi_v=-p.
	\label{eq:true_force_identification}
\end{equation}
Substitution into \eqref{eq:generalized_yield} gives the yield condition in
terms of pressure and equivalent shear stress,
\begin{equation*}
	f(p,q)=q-Mp\le0,
	\qquad
	M=\mu+\beta.
\end{equation*}
The generalized forces therefore coincide with the physical stress invariants
through \(\chi_q=q\) and \(\chi_v=-p\).  After this substitution, the yield
surface has volumetric slope \(M=\mu+\beta\), whereas the plastic rate in
\eqref{eq:generalized_loading_conditions} has volumetric slope \(\beta\).
Written directly in terms of pressure and equivalent shear stress, the plastic
rate therefore follows from the potential
\begin{equation}
	g(p,q)=q-\beta p,
	\qquad
	\dot{\bs\varepsilon}^{p}
	=\dot\lambda\frac{\partial g}{\partial\bs\sigma}.
	\label{eq:plastic_potential}
\end{equation}
Here the pressure, deviatoric stress, and equivalent shear stress are
\begin{equation*}
	p=-\frac{1}{3}\tr\bs\sigma,
	\qquad
	\mbf s=\dev\bs\sigma,
	\qquad
	q=\sqrt{\frac{3}{2}}\lVert\mbf s\rVert.
\end{equation*}
For \(q>0\), substitution into \eqref{eq:plastic_potential} gives
\begin{equation*}
	\dot{\bs\varepsilon}^{p}
	=\dot\lambda\left(
	\frac{3}{2}\frac{\mbf s}{q}
	+\frac{\beta}{3}\mbf I
	\right),
	\qquad
	\dot\gamma=\dot\lambda,
	\qquad
	\tr\dot{\bs\varepsilon}^{p}=\beta\dot\lambda.
\end{equation*}
Normality to \(f=0\) would give the volumetric coefficient \(M\), whereas the
actual coefficient is \(\beta\).  Because \(M=\mu+\beta\), the flow is
non-associated whenever \(\mu>0\).

The dissipation now follows directly from the active surface and the
dilatancy constraint,
\begin{align}
	\mathcal D^p
	 & =q\dot\gamma-p\dot\zeta^p
	\notag                              \\
	 & =\left(M-\beta\right)p\dot\gamma
	\notag                              \\
	 & =\mu p\dot\gamma
	\ge0.
	\label{eq:drucker_prager_dissipation}
\end{align}
Thus the frictional strength separates into the dilation slope \(\beta\) and
the part \(\mu\) that produces dissipation.  For \(\mu=0\), the flow is
associated and this idealized rate-independent mechanism dissipates no power.
Collins and Houlsby obtain the same separation in their stress-dependent
dissipation construction \cite[eq.~(5.14)]{collins1997}.

For \(p\ge0\), \(\dot\gamma\ge0\), and constants \(\mu,\beta\ge0\), this
construction gives the Drucker--Prager model exactly.  The zero-rate work
inequality gives the admissible force region in
\eqref{eq:generalized_yield}.  Substituting the physical forces from
\eqref{eq:true_force_identification} gives the yield slope
\(M=\mu+\beta\), while associated flow in the two-force description gives the
dilation slope \(\beta\).  Substitution into the plastic power gives
\eqref{eq:drucker_prager_dissipation}.  The yield and flow slopes differ when
\(\mu=M-\beta>0\), which is precisely the non-associated case.

The exact equivalence applies to an isotropic material when the true-plastic
rate remains aligned with its driving stress, loading is monotone with
\(p\ge0\) and \(\dot\gamma\ge0\), \(\mu\) and \(\beta\) are constant, and
cohesion and back stress are zero.  Cohesion and hardening can be introduced by
replacing \(\mu p\) with a nonnegative state-dependent strength and prescribing
its evolution.  Changes in loading direction, corners, cyclic reversal, and
anisotropy require a rate potential that specifies the tensor direction and
rules that determine which yield condition is active.  These requirements
distinguish the coupled model from the model with separate yield conditions in
Section~\ref{sec:two_mechanisms}.

\subsection{Relation to prior dissipation formulations}
\label{sec:prior_dissipation}

Previous work has already derived non-associated Drucker--Prager flow from
dissipation principles.  Collins and Houlsby impose the same linear dilation
constraint and stress-dependent dissipation used in
\eqref{eq:dissipation_potential}, obtaining the same separation
\(M=\mu+\beta\) \cite[Sec.~5]{collins1997}.  Equations
\eqref{eq:zero_rate_force_condition}--\eqref{eq:true_force_identification}
state the admissible forces at zero plastic rate and then substitute pressure
for the volumetric driving force.  They do not define a different small-strain
constitutive law.

The term \emph{coupled mechanisms} used here refers to a common activation
condition and plastic multiplier for the isochoric and volumetric mechanisms.
It is distinct from Collins and Houlsby's elastic--plastic coupling through the
thermodynamic potential \cite[Sec.~3]{collins1997}.

Srinivasa instead lets the dissipation function determine both the yield
surface and the non-associated flow direction, without imposing a separate
dilation constraint \cite{srinivasa2010}.  Ulloa et al. derive
non-associated flow from stored energy and a state-dependent dissipation
potential; their approach also includes hardening and more general
non-associated models \cite{ulloa2021}.  The present construction is narrower
than those approaches.  It assigns the shear and volume rates to distinct
true-plastic and distension factors, makes the closure requirement explicit,
compares independent activation with a shared-multiplier response, and carries
that mechanism-resolved interpretation to finite deformation.  Under the Flory
closure, these factors are the conventional isochoric and volumetric plastic
factors rather than additional degrees of freedom.

\subsection{Proportional loading comparison}
\label{sec:loading_comparison}

A small-strain proportional path shows how the two models respond.  Let the
elastic shear and bulk moduli be \(G>0\) and
\(K_e>0\), and impose the coaxial strain rate
\begin{equation*}
	\dot{\bs\varepsilon}
	=\dot\varepsilon_s\mbf N
	-\frac{\dot\varepsilon_v}{3}\mbf I,
	\qquad
	\dot\varepsilon_s>0,
	\qquad
	\dot\varepsilon_v>0,
\end{equation*}
where \(\dot\varepsilon_s\) is the equivalent deviatoric strain rate and
\(\dot\varepsilon_v=-\tr\dot{\bs\varepsilon}\) is the
compression-positive volumetric strain rate.  Take constant independent
strengths \(Y_0\) and \(K_0\), and start at their common corner
\(q=Y_0\), \(p=K_0\).  Choosing \(Y_0=MK_0\) places the coupled model at the
same stress point on its Drucker--Prager surface.

For the independent model, associated distension flow at \(p=K_0>0\) gives
\(\dot\zeta^p=-\dot\lambda_v\).  The elastic stress rates are
\begin{equation*}
	\dot q=3G(\dot\varepsilon_s-\dot\lambda_d),
	\qquad
	\dot p=K_e(\dot\varepsilon_v-\dot\lambda_v).
\end{equation*}
Consistency of both active surfaces therefore gives
\begin{equation*}
	\dot\lambda_d=\dot\varepsilon_s,
	\qquad
	\dot\lambda_v=\dot\varepsilon_v,
	\qquad
	\dot q=0,
	\qquad
	\dot p=0,
	\qquad
	\dot\zeta^p=-\dot\varepsilon_v<0.
\end{equation*}
The independent response remains at the corner and accommodates the imposed
volume reduction by plastic compaction.

For the coupled model, \(\dot\zeta^p=\beta\dot\lambda\), and the stress rates
under the same imposed path are
\begin{equation*}
	\dot q=3G(\dot\varepsilon_s-\dot\lambda),
	\qquad
	\dot p=K_e(\dot\varepsilon_v+\beta\dot\lambda).
\end{equation*}
Consistency of \(q-Mp=0\) gives
\begin{equation*}
	\dot\lambda
	=\frac{3G\dot\varepsilon_s-MK_e\dot\varepsilon_v}
	{3G+MK_e\beta}.
\end{equation*}
For \(\beta>0\) and
\(3G\dot\varepsilon_s>MK_e\dot\varepsilon_v\), the multiplier is positive,
\(\dot\zeta^p=\beta\dot\lambda>0\), and
\(\dot q=M\dot p>0\).  Thus the coupled response moves along the
Drucker--Prager surface while dilating, whereas the independent response under
the same loading remains at its two-surface corner while compacting.

This comparison also identifies the material behavior represented by each
construction.  The coupled model is appropriate when shear rearrangement causes
plastic volume expansion, as in the shear dilation of a dense frictional
granular material.  The independent model is appropriate when volume change has
its own activation process, as in pore collapse or grain crushing.  A material
that exhibits both shear dilation and pressure-driven compaction would generally
require a coupled shear-dilation mechanism together with a separate compaction
mechanism.

\section{Finite-deformation extension}
\label{sec:finite_lift}

The same stresses and internal rates that form the plastic power in
\eqref{eq:plastic_dissipation} also define the finite-deformation model.  This section adopts
the Flory closure \eqref{eq:flory_closure}; thus \(a=J\),
\(\bar J=1\), and \(a^e=J/a^p\) are determined from the imposed deformation
and the plastic state.  For this specialization, the model neglects energetic
hardening.  Active flow is restricted to \(q>0\), and the normal is
\begin{equation}
	q=\sqrt{\frac{3}{2}}\lVert\mbf M\rVert,
	\qquad
	\mbf N=\frac{3}{2}\frac{\mbf M}{q}.
	\label{eq:finite_equivalent_stress}
\end{equation}
The two internal mechanisms then evolve according to
\begin{subequations}
	\label{eq:finite_flow_rules}
	\begin{align}
		\dot{\bar{\mbf F}}^{p}(\bar{\mbf F}^{p})^{-1}
		 & =\dot\lambda\mbf N,
		\label{eq:finite_true_flow}      \\
		\dot\zeta^p
		 & =\frac{\dot a^p}{a^p}, \notag \\
		\dot\zeta^p
		 & =\beta\dot\lambda.
		\label{eq:finite_distension_flow}
	\end{align}
\end{subequations}
Because \(\tr\mbf N=0\), the first equation preserves
\(\det\bar{\mbf F}^{p}=1\), while the second evolves plastic distension.
The normalization in \eqref{eq:finite_equivalent_stress} gives
\(\dot\gamma=\dot\lambda\) during plastic loading.
Together, the two equations update the factors in
\begin{equation*}
	\mbf F
	=\left(a^e a^p\right)^{1/3}
	\bar{\mbf F}^{e}\bar{\mbf F}^{p}.
\end{equation*}

The yield condition and loading rules are
\begin{equation}
	f(\mbf M,p)=q-(\mu+\beta)p\le0,
	\qquad
	\dot\lambda\ge0,
	\qquad
	\dot\lambda f=0,
	\qquad
	\dot f=0\ \text{when}\ \dot\lambda>0.
	\label{eq:finite_loading_conditions}
\end{equation}
Substituting \eqref{eq:finite_flow_rules} and
\eqref{eq:finite_loading_conditions} into the finite-deformation plastic power
gives
\begin{align*}
	\mathcal D^p
	 & =\mbf M:
	\left[
		\dot{\bar{\mbf F}}^{p}(\bar{\mbf F}^{p})^{-1}
		\right]
	-p\frac{\dot a^p}{a^p}
	\\
	 & =q\dot\lambda-p\beta\dot\lambda
	\\
	 & =\mu p\dot\lambda
	\ge0.
\end{align*}
The dissipation proof pairs each internal rate with the stress that does work
through it.  The finite-deformation yield surface is therefore a
Drucker--Prager-type condition in \((p,q)\), with \(p\) formed from the
Kirchhoff stress and \(q\) from the Mandel-type true-plastic stress.  This
surface reduces to the conventional Cauchy-stress form at small strain.  At
finite strain, a particular Drucker--Prager implementation must express its
chosen stress measure and stress update in terms of \(\mbf M\) and \(p\) before
the two formulations can be compared.

The same associated-flow argument applies at finite deformation when pressure
is held fixed.  With the generalized thermodynamic forces \(\chi_q\) and
\(\chi_v\) conjugate to \(\dot\gamma\) and \(\dot\zeta^p\), the
admissible-force boundary, the physical stress invariants, and the plastic
rates are
\begin{equation*}
	\chi_q+\beta\chi_v-\mu p=0,
	\qquad
	\chi_q=q,
	\qquad
	\chi_v=-p,
	\qquad
	(\dot\gamma,\dot\zeta^p)=\dot\lambda(1,\beta).
\end{equation*}
Substituting \(\chi_q=q\) and \(\chi_v=-p\) gives
\eqref{eq:finite_loading_conditions}, while the plastic rate retains the
dilation slope \(\beta\).  Thus the model is associated in the two driving
forces but non-associated when written in terms of \((p,q)\) whenever
\(\mu>0\).

\section{Discussion}
\label{sec:discussion}

The split \(\mbf F=\mbf A\bar{\mbf F}\) gives the two plastic rates distinct
physical roles.  True-plastic flow rearranges material without changing its
volume, whereas plastic distension changes the volume accommodated by that
material.  The free-energy inequality shows how each mechanism contributes to
plastic power, and any constitutive model must keep their sum nonnegative.  When a
microstructural volume variable closes the general split, these rates can
represent independent mechanisms.  Under the Flory closure used here, they
are the isochoric and volumetric parts of conventional plastic flow.

The independent and coupled models give different responses.  With separate
yield conditions, either mechanism can flow alone, or both can flow with
independent multipliers.  A one-surface model has only one loading condition
and one multiplier, so it cannot reproduce all three cases without an
additional coupling rule.  The Drucker--Prager model provides that rule through
a common multiplier, a dilatancy constraint, and stress-dependent dissipation.

The decomposition also identifies how a dilating frictional solid dissipates
energy.  Under compression, plastic distension contributes negative power
during dilation, while the true-plastic mechanism contributes sufficient
positive power to make their sum \(\mu p\dot\gamma\).  If each contribution
were required to be nonnegative on its own, the model would allow compaction
but not dilation.  Dilation is possible because the positive isochoric
contribution offsets the negative volumetric contribution.

The present theory applies locally to isotropic, cohesionless materials under
monotone compression with \(q>0\).  Elasticity, hardening or softening,
reversal, and anisotropy
require additional constitutive choices that preserve the power in
\eqref{eq:plastic_dissipation}.  Under the adopted Flory closure, the displayed
stress measures and rate equations fully specify the plastic response
considered here.

\section{Conclusions}
\label{sec:conclusions}

This paper gives apparent non-associated plasticity a mechanism-resolved
finite-deformation formulation based on the split
\(\mbf F=\mbf A\bar{\mbf F}\).  The free-energy inequality separates
true-plastic deformation from plastic distension and identifies their
conjugate forces as a deviatoric Mandel-type stress and the negative of the
compression-positive mean stress.  Each independent mechanism follows an
associated flow rule, but their combined plastic strain rate need not be normal
to either yield surface.

For the stated isotropic, cohesionless, monotone compression regime, coupling
the mechanisms through a state-dependent dissipation potential gives an exact
Drucker--Prager specialization.  The construction yields the yield slope
\(M=\mu+\beta\), the dilation slope \(\beta\), and the nonnegative power
\(\mu p\dot\gamma\).  In terms of the two forces that drive the internal
mechanisms, the flow is associated.  Substituting pressure and shear stress for
those forces gives the plastic potential \(g=q-\beta p\) and the familiar
non-associated flow rule.  Under the Flory closure, the same division of
plastic power carries to the conventional finite-deformation volumetric and
isochoric plastic factors.  The small-strain law agrees with the
Collins--Houlsby generalized-stress construction; the additional content is the
identification of its kinematic carriers and of the activation and closure
needed to interpret them as distinct physical mechanisms.

\section*{Acknowledgements}

The author thanks Professor D. Nicolas Espinoza.  This work was funded by the
U.S. Department of Energy under awards DE-FE0032349, ``Hydrogen Storage in Salt
Caverns in the Permian Basin: Seal Integrity Evaluation and Field Test,'' and
DE-AR0001864283, assigned control number 2784-1755, ``Cyclic Injection for
Commercial Seismic-Safe Geologic H$_2$ Production.''

\section*{Declaration of generative AI and AI-assisted technologies}

From 8 May 2026 through the latest covered commit on 19 August 2026, the author used OpenAI Codex (GPT-5.6 Terra and GPT-5.6 Sol) and GitHub Copilot (DeepSeek-V4-Flash) within an author-directed, version-controlled manuscript-development workflow. AI assistance was used to propose and critique prose, examine notation and equation references, check derivations for consistency, identify literature and bibliography issues, and check the LaTeX output. The author determined the scientific questions, theoretical structure, assumptions, mathematical arguments, physical interpretations, and final wording. Every accepted change was reviewed by the author, and equations and citations were checked against manuscript source and primary literature as applicable. AI output was not treated as a scholarly source or credited with authorship. The author takes full responsibility for the accuracy, originality, and integrity of the work.
The repository at \url{https://github.com/johntfoster/nonassociated_plasticity} is available for inspection of the manuscript's changes and development history; it does not include the specific prompts used.

\bibliographystyle{plain}
\bibliography{references}

\end{document}
